\documentclass[../main.tex]{subfiles} % 使用 subfiles 文档类，指定主文档
\graphicspath{{\subfix{../image/}}} % 指定图片目录，后续可以直接使用图片文件名。

% 例如:
% \begin{figure}[H]
% \centering
% \includegraphics[width=0.3\textwidth]{图.png}
% \caption{}
% \label{figure:图}
% \end{figure}
% 注意:上述\label{}一定要放在\caption{}之后,否则引用图片序号会只会显示??.

\begin{document}
% SOURCE p=100; printed=85; object=ch03-title; region=top
\chapter{符号测度与微分}\label{chapter:ch03-signed-measures-differentiation}
% SOURCE p=100; printed=85; object=ch03-intro; region=middle
本章的主要主题，是对定义在同一个 $\sigma$-代数上的测度 $\nu$ 关于另一个测度 $\mu$ 进行微分。我们先在抽象情形下研究这一问题，再对 $\mu$ 为 $\mathbb{R}^n$ 上 Lebesgue 测度的特殊情形作进一步分析。在后一种情形下，当 $n=1$ 时，我们与经典实变量理论建立联系，并得到 Lebesgue 积分的微积分基本定理。

首先，要把测度的概念推广到允许取负值乃至复值的情形。这有三个理由。第一，在应用中，这种更一般的测度常会出现；例如，电荷既可以为正，也可以为负。第二，它们是微分理论的自然研究对象。第三，复测度具有重要的泛函分析意义，我们将在\cref{chapter:ch07-radon-measures}中讨论。
\subfile{Chapter-03/section-01.tex}
\subfile{Chapter-03/section-02.tex}
\subfile{Chapter-03/section-03.tex}
\subfile{Chapter-03/section-04.tex}
\subfile{Chapter-03/section-05.tex}
\subfile{Chapter-03/section-06.tex}
\end{document}
